% ---------------------------------------------------------------------------
% Section: Phase 19 -- the engine update and LittleBit
% ---------------------------------------------------------------------------
\section{Phase 19: The Engine Update and LittleBit}
\label{sec:phase19}

Phase~19 (5 November 2026) replaces the vendored Smiley Core (revision 7cb6028) by the owner's new one (06c2b88), re-applies Nafs-Aql's additions on it,
integrates LittleBit~\cite{littlebit2025}, and re-validates the system. The decision record is \texttt{docs/decisions/013-smiley-core-littlebit.md} and
the research note \texttt{docs/research/smiley\_core\_update.md}.

\subsection{What the new engine brings}

Grouped online-softmax attention, an int8 key--value cache rotated by a Hadamard transform, vectorised math, and sampling whose randomness is keyed by
position, so a drafted token is accepted exactly when sequential sampling would have emitted it: speculation becomes lossless at any temperature, not
only under argmax. LittleBit approximates a weight matrix as a product of sign matrices and small scale vectors, $\hat W=\mathrm{diag}(h)\,U_{\pm}\,
\mathrm{diag}(\ell)\,V_{\pm}^{\top}\,\mathrm{diag}(g)$, summed over stages, stored beside the model as a sidecar file.

\subsection{The merge}

Nafs-Aql's vendored copy differs from upstream by a batching engine, a prefix store, a continuation prefix, log-probabilities and two sampler cuts, so
the update was a three-way merge (base 7cb6028) after normalising both sides with the same formatter, which reduced 44 textual conflicts to 14. The
additions were then re-applied on the new engine: a cache snapshot that records whether the cache is rotated; the per-row forward pass and attention
without touching upstream's attention kernel (one row per user, so the bit-identity property of Phase~18 holds); top-$k$ and min-$p$ on the
position-keyed sampler, in llama.cpp's order, with upstream's fast path when both are off. All 370 tests of the workspace pass in release mode,
including the batching bit-identity tests, and formatting and the linter are clean.

\subsection{Two defects the integration exposed}

\emph{A crash.} Served through the batching engine, a LittleBit model crashed on its first request: the single-session path sizes the sidecar's scratch
buffers and the engine's constructor did not, so the first sidecar matmul wrote through an empty buffer. Fixed, and visible only because the engine
was exercised with the new format.

\emph{A fourfold slowdown with no kernel change.} After the merge Qwen3-1.7B decoded at 4.4--5.0 tokens per second, against 17--20 for the previous
commit and 16--18 for upstream's own build, interleaved on the same machine. The dot kernels were byte-identical after formatting, so a code change
could not be the cause directly. Bisecting by swapping files between a scratch copy of upstream and the crate showed that merely adding the batching
engine (never called by the benchmark) was enough; a kernel micro-benchmark isolated it to the small-batch dot kernels (Q4\_K, one activation, one
thread: 679\,ns per row against 96); disassembly showed a call to \texttt{core::array::try\_from\_fn} \emph{inside} the kernel, once per super-block,
in the slow build only. \texttt{std::array::from\_fn} in two helpers marked \texttt{inline(always)} is inlined or not depending on how much unrelated
code the crate contains. Spelling the eight-element arrays out restores 96\,ns per row and 15--20 tokens per second. The lesson is general: a kernel
whose speed depends on an optimiser decision about code elsewhere has to be guarded by a benchmark, and \texttt{kbench} is now part of the crate for that.

\begin{table}[h]
\centering
\caption{Decode speed of Qwen3-1.7B Q4\_K\_M (tokens per second, median of five interleaved rounds of 96 tokens, three threads, one slot, shared host).}
\label{tab:phase19speed}
\begin{tabular}{lr}
\toprule
Engine & tok/s \\
\midrule
Previous (7cb6028 + Phase 18) & 18.7 \\
New (06c2b88), after the fix & 19.2 \\
New, before the fix & 4.4--5.0 \\
New with a LittleBit sidecar & 20.1 \\
\bottomrule
\end{tabular}
\end{table}

\subsection{LittleBit: integrated, tested, not useful yet}

Qwen3-1.7B was converted at 1.0 bit per weight with the engine's post-training (Dual-SVID) converter: 196 projections, 172.7\,MB of factors for 845.7\,MB
of weights ($4.9\times$), 0.980 bits per weight, 13 minutes on four cores. The mean relative reconstruction error of the weights is 0.752 (Q4\_K itself
is about 0.1). On twelve exact-answer prompts, greedy, interleaved on the same engine, \textbf{the dense model answered 9 correctly and the LittleBit
model 0}; its text is incoherent. Decode was only 5\,\% faster (20.1 against 19.2\,tok/s), because at this size the matmul is not the whole token. Upstream
reports the same on a smaller model and attributes the paper's quality to distillation-based quantisation-aware training, which needs a GPU. The
integration is therefore opt-in (\texttt{--littlebit}, \texttt{LITTLEBIT=1}); a missing or corrupt sidecar refuses to start, so a stale file beside a
model cannot switch the weights silently, and the weights in use (dense or LittleBit) are reported in every response, in the backend health check and
on the Status page. It is ready for trained factors; it is not a result.

\subsection{Is anything else different}

The old and the new engine give the same answers on the exact-answer prompts (9 of 12 on both, the same three misses). One of those misses, a
time-of-day sum (3:40\,pm plus 2\,h 35\,min answered 5:45\,pm, the correct answer being 6:15\,pm), was first seen in the web interface right after the merge and
looked like a regression; the comparison shows it is the 1.7B model's own error, which the gate's confidence signal does not catch because both quick
samples repeat it. Re-validation on the final code: contract tests, 26 of 26 security probes (0 open high-severity findings), 59 of 59 Rooms checks, 65 of 65
end-to-end checks, 53 of 53 harness checks and 9 of 9 acceptance tests pass.
